% GENERATED FILE. Edit canonical lessons, then run npm run generate:books.
% canonical-source-hash: fnv1a64:06f2b435c1720a7e
% canonical-lessons: FR-C128-musique, FR-C128-sport, FR-C128-jeu, FR-C128-film, FR-C128-chanson

\chapter{Music, Sport, Game, Film, Song}
\label{ch:fr-music-sport-game-film-song}
\hlchaptermodality{128}

\begin{chapteropening}
\textbf{By the end of this chapter:} \emph{I can name music, sport, a game, a film, a song.}

Complete the last lesson of chapter 128: I can name music, sport, a game, a film, a song.
\end{chapteropening}

\section[la musique]{la musique — music}

\label{lesson:FR-C128-musique}



\begin{quote}
Before the new one: say the French for an idea, then the French for a problem.

Say \emph{j'aime bien la musique}. \textbf{Aimer bien} is to like, softer than \textbf{aimer}.
\end{quote}

\subsection*{You'll want to know: la musique}
\textbf{la musique} — \textquotedblleft{}music\textquotedblright{}.

\begin{cousinweb}
\textbf{musique} $\leftarrow$ Latin \emph{musica}, from Greek \emph{mousike}, \textquotedblleft{}the art of the Muses\textquotedblright{} — a learned borrowing. (EtymDB 2.1)
\end{cousinweb}

\begin{grammarlens}[title={What you've built}]
One more word for this chapter.
\end{grammarlens}

\subsection*{Your turn}
\begin{itemize}
\raggedright
  \item \emph{Say it:} \emph{la musique}
  \item \emph{Say it:} \emph{la musique}, once more
  \item \emph{Say it:} say \emph{la musique}
  \item \emph{Recall:} say the French for an idea, then the French for a problem, then say \emph{la musique} again
\end{itemize}

\begin{tcolorbox}[breakable,colback=teal!4,colframe=teal!35!black,title={Before you move on}]
What does la musique mean? (\textquotedblleft{}Music\textquotedblright{}.) Say it once more.
\end{tcolorbox}
\section[le sport]{le sport — sport}

\label{lesson:FR-C128-sport}



\begin{quote}
Before the new one: say the French for a problem, then the French for music.

Sport on a Sunday: say \emph{dimanche}.
\end{quote}

\subsection*{You'll want to know: le sport}
\textbf{le sport} — \textquotedblleft{}sport\textquotedblright{}.

\begin{cousinweb}
\textbf{sport} $\leftarrow$ English \emph{sport} — borrowed in the 19th century. English had shortened it from Old French \emph{desport}, \textquotedblleft{}a pastime\textquotedblright{}. (EtymDB 2.1)
\end{cousinweb}

\begin{grammarlens}[title={What you've built}]
One more word for this chapter.
\end{grammarlens}

\subsection*{Your turn}
\begin{itemize}
\raggedright
  \item \emph{Say it:} \emph{le sport}
  \item \emph{Say it:} \emph{le sport}, once more
  \item \emph{Say it:} say \emph{le sport}
  \item \emph{Recall:} say the French for a problem, then the French for music, then say \emph{le sport} again
\end{itemize}

\begin{tcolorbox}[breakable,colback=teal!4,colframe=teal!35!black,title={Before you move on}]
What does le sport mean? (\textquotedblleft{}Sport\textquotedblright{}.) Say it once more.
\end{tcolorbox}
\section[le jeu]{le jeu — a game}

\label{lesson:FR-C128-jeu}



\begin{quote}
Before the new one: say the French for music, then the French for sport.
\end{quote}

\subsection*{You'll want to know: le jeu}
\textbf{le jeu} — \textquotedblleft{}a game\textquotedblright{}. Plural \textbf{les jeux}, with an x.

\begin{cousinweb}
\textbf{jeu} $\leftarrow$ Old French \emph{geu} $\leftarrow$ Latin \emph{iocus}, \textquotedblleft{}a joke, a jest\textquotedblright{} — inherited. It is the noun beside \emph{jouer}. (EtymDB 2.1)
\end{cousinweb}

\begin{grammarlens}[title={What you've built}]
One more word for this chapter.
\end{grammarlens}

\subsection*{Your turn}
\begin{itemize}
\raggedright
  \item \emph{Say it:} \emph{le jeu}
  \item \emph{Say it:} \emph{le jeu}, once more
  \item \emph{Say it:} say \emph{le jeu}
  \item \emph{Recall:} say the French for music, then the French for sport, then say \emph{le jeu} again
\end{itemize}

\begin{tcolorbox}[breakable,colback=teal!4,colframe=teal!35!black,title={Before you move on}]
What does le jeu mean? (\textquotedblleft{}A game\textquotedblright{}.) Say it once more.
\end{tcolorbox}
\section[le film]{le film — a film}

\label{lesson:FR-C128-film}



\begin{quote}
Before the new one: say the French for sport, then the French for a game.
\end{quote}

\subsection*{You'll want to know: le film}
\textbf{le film} — \textquotedblleft{}a film\textquotedblright{}.

\begin{cousinweb}
\textbf{film} $\leftarrow$ English \emph{film}, \textquotedblleft{}a thin skin, a membrane\textquotedblright{}, then the coated strip a camera uses — borrowed. (EtymDB 2.1)
\end{cousinweb}

\begin{grammarlens}[title={What you've built}]
One more word for this chapter.
\end{grammarlens}

\subsection*{Your turn}
\begin{itemize}
\raggedright
  \item \emph{Say it:} \emph{le film}
  \item \emph{Say it:} \emph{le film}, once more
  \item \emph{Say it:} say \emph{le film}
  \item \emph{Recall:} say the French for sport, then the French for a game, then say \emph{le film} again
\end{itemize}

\begin{tcolorbox}[breakable,colback=teal!4,colframe=teal!35!black,title={Before you move on}]
What does le film mean? (\textquotedblleft{}A film\textquotedblright{}.) Say it once more.
\end{tcolorbox}
\section[la chanson]{la chanson — a song}

\label{lesson:FR-C128-chanson}



\begin{quote}
Before the new one: say the French for a game, then the French for a film.
\end{quote}

\subsection*{You'll want to know: la chanson}
\textbf{la chanson} — \textquotedblleft{}a song\textquotedblright{}.

\begin{cousinweb}
\textbf{chanson} $\leftarrow$ Old French \emph{chançon} $\leftarrow$ Latin \emph{cantio}, \emph{cantionis}, \textquotedblleft{}a song\textquotedblright{}, from \emph{canere}, \textquotedblleft{}to sing\textquotedblright{} (EtymDB 2.1) — inherited: \emph{ca-} became \emph{cha-}.
\end{cousinweb}

\begin{grammarlens}[title={What you've built}]
One more word for this chapter.
\end{grammarlens}

\subsection*{Your turn}
\begin{itemize}
\raggedright
  \item \emph{Say it:} \emph{la chanson}
  \item \emph{Say it:} \emph{la chanson}, once more
  \item \emph{Say it:} say \emph{la chanson}
  \item \emph{Recall:} say the French for a game, then the French for a film, then say \emph{la chanson} again
\end{itemize}

\begin{tcolorbox}[breakable,colback=teal!4,colframe=teal!35!black,title={Before you move on}]
What does la chanson mean? (\textquotedblleft{}A song\textquotedblright{}.) Say it once more.
\end{tcolorbox}
